\chapter{近世代数：正交群怎样作用在二次型上}
\label{chap:abstract}

把两个图形比较一次，只需找到一个保距变换。若要研究\emph{所有}可能的转动、反射和摆放方式，就需要能把这些变换组织起来的语言。群描述变换的复合；群作用描述这些变换怎样改变对象。本章只考虑固定原点的变换，把二次型的系数矩阵作为对象，看同一个中心椭圆可以有多少种摆法，以及哪些变换使它保持原状。

\section{旋转与反射的复合}
一个\textbf{群}具有封闭的乘法，乘法满足结合律，有单位元，且每个元素都有逆元。所有满足 \(U^TU=I\) 的 \(n\) 阶实矩阵在矩阵乘法下组成正交群 \(O(n)\)。因为
\[
 (UV)^T(UV)=V^TU^TUV=I,\qquad U^{-1}=U^T,
\]
两个正交变换可以连续做，也可以各自撤销。由 \((\det U)^2=1\) 可知 \(\det U=\pm1\)。其中行列式为 \(1\) 的矩阵构成特殊正交群 \(SO(n)\)。

在平面上，设 \(R_\theta\) 是逆时针旋转 \(\theta\) 的矩阵，设
\[
 R_\theta=\begin{pmatrix}\cos\theta&-\sin\theta\\
 \sin\theta&\cos\theta\end{pmatrix},\qquad
 S=\begin{pmatrix}1&0\\0&-1\end{pmatrix}.
\]
若 \(U\in O(2)\) 的第一列为 \((\cos\theta,\sin\theta)^T\)，第二列只能是 \((-\sin\theta,\cos\theta)^T\) 或它的相反向量。故 \(U\) 或为 \(R_\theta\)，或为 \(R_\theta S\)。前者是旋转；后者是先关于横轴反射，再旋转。

\begin{example}{两次反射，为什么转过两倍角}
先关于横轴反射，再关于与横轴成 \(30^\circ\) 的直线反射。求合成变换，并求 \((2,0)^T\) 的像。若交换反射次序，结果怎样？
\end{example}
\begin{solution}
先约定角度的正方向：\(\beta\) 从正横轴逆时针量起，表示\textbf{反射直线的方向角}，不是点被旋转的角。直线没有指定朝向，所以 \(\beta\) 与 \(\beta+180^\circ\) 表示同一条反射线；反射矩阵也应相同。以下固定坐标轴、主动移动点，并使用列向量。\textbf{矩阵乘积中，右边的变换先做。}

\textbf{第一步：写出关于斜线的反射。}先把点旋转 \(-\beta\)，使反射直线转到横轴；关于横轴反射；最后旋转 \(\beta\) 回到原来的位置。因此
\[
 F_\beta=R_\beta S R_{-\beta}
 =\begin{pmatrix}
 \cos 2\beta&\sin 2\beta\\
 \sin 2\beta&-\cos 2\beta
 \end{pmatrix}.
\]
矩阵中的 \(2\beta\) 并不表示 \(F_\beta\) 是旋转：它的行列式为 \(-1\)，而且整条反射直线上的点都不动。若一个非零向量原方向角为 \(\varphi\)，反射后方向角为 \(2\beta-\varphi\)。

\textbf{第二步：按题目次序复合。}横轴反射把方向角 \(\varphi\) 变为 \(-\varphi\)，再作斜线反射得到 \(2\beta+\varphi\)，所以每个方向都统一增加 \(2\beta\)。矩阵计算也给出
\[
 F_\beta S
 =\begin{pmatrix}
 \cos 2\beta&-\sin 2\beta\\
 \sin 2\beta&\cos 2\beta
 \end{pmatrix}
 =R_{2\beta}.
\]
交换次序则为
\[
 SF_\beta
 =\begin{pmatrix}
 \cos 2\beta&\sin 2\beta\\
 -\sin 2\beta&\cos 2\beta
 \end{pmatrix}
 =R_{-2\beta}.
\]
这两个合成变换互为逆变换。因为每次反射都做两遍就回到变换前的位置，即 \(S^2=F_\beta^2=I\)，所以撤销 \(F_\beta S\) 必须按相反次序做 \(SF_\beta\)。

\textbf{第三步：代入 \(\beta=30^\circ\)。}此时
\[
 F_{30^\circ}=\frac12\begin{pmatrix}1&\sqrt3\\\sqrt3&-1\end{pmatrix},
 \quad F_{30^\circ}S=R_{60^\circ},
 \quad SF_{30^\circ}=R_{-60^\circ}.
\]
对题目中的具体点，两种次序的每一步分别为
\begin{align*}
 (2,0)^T&\xrightarrow{\ S\ }(2,0)^T
          \xrightarrow{\ F_{30^\circ}\ }(1,\sqrt3)^T,\\
 (2,0)^T&\xrightarrow{\ F_{30^\circ}\ }(1,\sqrt3)^T
          \xrightarrow{\ S\ }(1,-\sqrt3)^T.
\end{align*}
两像的长度都是 \(\sqrt{1+3}=2\)，但点的位置不同。反向合成把正向的像还原：\(R_{-60^\circ}(1,\sqrt3)^T=(2,0)^T\)。

\textbf{第四步：核对保距。}直接相乘可得 \(F_\beta^TF_\beta=I\) 及 \(R_\theta^TR_\theta=I\)。所以对任意两点 \(p,q\)，各次反射及两个合成变换 \(T\) 都满足
\[
 \|T(p)-T(q)\|^2
 =(p-q)^TT^TT(p-q)=\|p-q\|^2.
\]
这是对所有点对成立的保距证明；具体点的长度检查只是其中一个例子。两次反射的合成虽保距，交换次序却通常改变结果。本题的 \(60^\circ\) 与 \(-60^\circ\) 旋转正好展示了这个区别。
\end{solution}

\section{轨道描述所有摆法，稳定子描述对称}
设 \(A\) 是实对称矩阵。定义左作用
\[
 U\cdot A=UAU^T \qquad (U\in O(n)).
\]
它确实满足 \(I\cdot A=A\) 和 \(U\cdot(V\cdot A)=(UV)\cdot A\)。这里的写法对应主动移动图形：若 \(x=Uz\)，则
\[
 x^T(U\cdot A)x=z^TAz.
\]
因此 \(U\) 把 \(A\) 的等值图形旋转或反射到 \(U\cdot A\) 的等值图形。所有 \(U\cdot A\) 组成 \(A\) 的\textbf{轨道}；满足 \(U\cdot A=A\) 的变换组成 \(A\) 的\textbf{稳定子}。谱定理给出每条轨道的对角代表；把特征值按大小排序后，这个对角代表唯一。

\begin{example}{只许转直角，椭圆有几种摆法}
令
\[
 A=\begin{pmatrix}5&2\\2&2\end{pmatrix},\qquad
 J=\begin{pmatrix}0&-1\\1&0\end{pmatrix}.
\]
在只允许 \(I,J,J^2,J^3\) 的情况下，求 \(A\) 的轨道和稳定子，并解释图形的含义。
\end{example}
\begin{solution}
\(J\) 是逆时针旋转 \(90^\circ\)，且 \(J^2=-I\)、\(J^4=I\)。直接计算
\[
 JAJ^T=\begin{pmatrix}2&-2\\-2&5\end{pmatrix}=:B.
\]
由于 \((-I)A(-I)^T=A\)，\(J^2\) 使矩阵回到 \(A\)，而 \(J^3=-J\) 给出 \(B\)。故轨道为 \(\{A,B\}\)，在这四种变换中的稳定子为 \(\{I,-I\}\)。
因此椭圆 \(5x^2+4xy+2y^2=1\) 旋转 \(90^\circ\) 得到不同摆法：例如 \((1/\sqrt5,0)\) 在原椭圆上，代入新方程左边却为 \(2/5\)。旋转 \(180^\circ\) 则回到原状。两个矩阵的迹都为 \(7\)，行列式都为 \(6\)，所以特征值同为 \(6,1\)；这也与正交共轭保谱相符。
\end{solution}

正定二次型的单位等值集足以确定其系数矩阵。事实上，若两个正定二次型 \(q_1,q_2\) 的单位等值集相同，任取 \(x\ne0\)，令 \(z=x/\sqrt{q_1(x)}\)，则 \(q_1(z)=q_2(z)=1\)。把缩放比例代回便得 \(q_1(x)=q_2(x)\)。对实对称矩阵，分别比较各平方项和交叉项的系数，就知道两个矩阵相同。因此，研究下一题的矩阵等式，正好等于研究相应椭圆的摆放。

\begin{example}{求出把标准椭圆移到斜椭圆的所有正交矩阵}
设 \(D=\operatorname{diag}(6,1)\)，并令
\[
 Q=\frac1{\sqrt5}\begin{pmatrix}2&-1\\1&2\end{pmatrix},
 \qquad A=QDQ^T.
\]
求所有满足 \(UDU^T=A\) 的 \(U\in O(2)\)。其中有几个是旋转？
\end{example}
\begin{solution}
先求 \(D\) 的稳定子。令 \(V=Q^TU\)，则 \(UDU^T=A\) 等价于 \(VDV^T=D\)，即 \(VD=DV\)。写 \(V=(v_{ij})\)，比较非对角元得
\[
 v_{12}=6v_{12},\qquad 6v_{21}=v_{21}.
\]
故 \(V\) 是对角矩阵。正交性又要求 \(v_{11},v_{22}\in\{1,-1\}\)。反过来，每个这样的 \(V\) 都固定 \(D\)，所以全部解为
\[
 U=Q\operatorname{diag}(\varepsilon_1,\varepsilon_2),
 \qquad \varepsilon_1,\varepsilon_2\in\{1,-1\}.
\]
这四个矩阵对应两个主轴方向可以各自反向。因为 \(\det Q=1\)，其中行列式为 \(1\) 的恰是 \(Q\) 和 \(-Q\)，两种旋转相差 \(180^\circ\)。其余两个是反射。椭圆的主轴方向由两条特征直线确定，因此改变特征向量的正负号不会改变主轴。
\end{solution}

\section{特征值重合以后，对称变多了}
\begin{example}{同一特征值的两个方向可以混合吗}
令 \(D_3=\operatorname{diag}(6,1,1)\)。求所有满足 \(UD_3U^T=D_3\) 的 \(U\in O(3)\)，并说明 \(6x^2+y^2+z^2=1\) 的相应对称。
\end{example}
\begin{solution}
由 \(UD_3=D_3U\) 得，对于第 \(i,j\) 个矩阵元 \(u_{ij}\)，有 \((d_j-d_i)u_{ij}=0\)，其中 \(d_1=6,d_2=d_3=1\)。所以第一行、第一列的非对角元必须为零，而右下角的 \(2\times2\) 块可以有非零的非对角元。再用正交性得到
\[
 U=\begin{pmatrix}\varepsilon&0\\0&W\end{pmatrix},
 \qquad \varepsilon=\pm1,\quad W\in O(2).
\]
每个这样的 \(U\) 都与 \(D_3\) 交换，故这就是全部稳定子。几何上，\(x\) 轴可以反向，\(yz\) 平面内可以任意旋转或反射。比如固定 \(x\)，把 \((y,z)\) 换成 \((-z,y)\)，表达式 \(y^2+z^2\) 保持不变。与上一题相比，特征值 \(1\) 的特征空间由一条直线变成一个平面，允许在其中任取一组正交单位基。
\end{solution}

\section{区分代数合同和欧氏全等}
一般可逆换元 \(x=Pz\) 给出 \(A\mapsto P^TAP\)，称为合同。当 \(P\) 正交时，\(P^T=P^{-1}\)，这个公式也可读作相似变换。一般合同保留二次型的惯性，即正、负、零平方项的个数；正交共轭还保留特征值的具体数值。这个差别决定了变换是否保留图形的轴长。

\begin{example}{惯性相同的两条双曲线一定全等吗}
比较 $2x^2-8y^2=1$ 与 $u^2-v^2=1$。找出一个合同换元，再判断两条双曲线是否全等。
\end{example}
\begin{solution}
取 $P=\operatorname{diag}(1/\sqrt2,1/\sqrt8)$，即 $x=u/\sqrt2,\ y=v/\sqrt8$。代入得到 $u^2-v^2=1$；两边都是一个正平方项和一个负平方项，惯性相同。但 $P$ 不保距离。要证明不存在其他等距变换，可比较两个连通分支之间的最小距离。第一条曲线的右支满足 $x\ge1/\sqrt2$，左支满足 $x\le-1/\sqrt2$。两支之间任意两点的距离至少为 $\sqrt2$，在 $(\pm1/\sqrt2,0)$ 处达到。第二条曲线的同一距离是 $2$，在 $(\pm1,0)$ 处达到。等距双射保持连通分支及分支间的最小距离，因此两条双曲线不全等。此例表明，惯性确定二次型能化成哪一类符号标准形，却不记录几何尺度。
\end{solution}

另一方面，仅凭某一个等值点集相同，未必能推出对应矩阵在同一正交轨道。例如 \(x^2-2y^2=0\) 与 \(2x^2-4y^2=0\) 都给出直线 \(x=\pm\sqrt2\,y\)，但对应特征值分别为 \(\{1,-2\}\) 和 \(\{2,-4\}\)。在零水平上整体乘非零常数不改变点集，却可能改变谱。

\section{有限群的平均：给定对称，反过来找内积}
若有限群 \(G\) 的每个元素 \(g\) 对应一个可逆矩阵 \(\rho(g)\)，且 \(\rho(gh)=\rho(g)\rho(h)\)、\(\rho(e)=I\)，便得到一个矩阵表示。原来的欧氏内积未必被它保留。把所有矩阵作用后的内积取平均，可得到新内积 \(x^THy\)，其中
\[
 H=\frac{1}{|G|}\sum_{g\in G}\rho(g)^T\rho(g).
\]
对非零 \(x\)，和式中的单位元项已给出正贡献，所以 \(H\) 正定。对任意 \(h\in G\)，右乘 \(h\) 会重新排列 \(G\) 中的元素，故
\[
 \rho(h)^T H\rho(h)
 =\frac1{|G|}\sum_{g\in G}\rho(gh)^T\rho(gh)=H.
\]
谱定理允许取 \(C=H^{1/2}\)，于是 \(C^TC=H\)。用坐标 \(z=Cx\) 表示时，各 \(\rho(h)\) 变成正交矩阵 \(C\rho(h)C^{-1}\)。

\begin{example}{把一个斜反射写成普通反射}
令 \(T=\begin{pmatrix}1&2\\0&-1\end{pmatrix}\)。验证 \(\{I,T\}\) 是群，求它平均得到的 \(H\)，并找坐标 \(z=Cx\)，使 \(T\) 在新坐标下是关于横轴的反射。
\end{example}
\begin{solution}
矩阵相乘得 \(T^2=I\)，且 \(T\ne I\)，故有一个二元群。原欧氏度量下，\(T(0,1)^T=(2,-1)^T\)，长度从 \(1\) 变为 \(\sqrt5\)，所以 \(T\) 不正交。平均得到
\[
 H=\frac{I+T^TT}{2}
  =\begin{pmatrix}1&1\\1&3\end{pmatrix}.
\]
取 \(C=\begin{pmatrix}1&1\\0&\sqrt2\end{pmatrix}\)，便有 \(C^TC=H\)。对 \(x=(a,b)^T\)，新坐标是 \(z=(a+b,\sqrt2 b)^T\)；而
\[
 CTx=C(a+2b,-b)^T=(a+b,-\sqrt2 b)^T.
\]
所以新坐标下恰为 \((z_1,z_2)\mapsto(z_1,-z_2)\)。在原坐标中，\(T\) 保持的新长度是 \(\sqrt{x^THx}=\sqrt{(a+b)^2+2b^2}\)。这说明一个有限群的表示可以在合适的内积下保距；改变内积与证明它原本就保持欧氏距离是两件事。
\end{solution}
